LLM judges are increasingly used to evaluate code changes, but the main
benchmark for them, CodeJudgeBench, contains only single-file programs whose
correctness is decided by hidden tests.  We ask whether judge reliability
survives the move to repository-level code changes whose defects pass the
tests.  We build \RJB{}: \NPairs{} (gold patch, defective variant) pairs from
\NInstances{} real pull requests created after 1~January~2025 in \NRepos{}
Python repositories.  Each variant is a template-guided AST mutation of the
real patch from one of six review-oriented defect classes, and it is admitted
only if every test that passes with the gold patch also passes with the
variant inside the task's published Docker environment; the pull requests'
tests caught \TestCatchRate{} of the candidates.  We evaluate six open-weight
judges (1.5--4B parameters, plus a thinking mode) with the CodeJudgeBench
pairwise protocol on \RJB{} and on a CodeJudgeBench control subset.  Raw
pairwise accuracy is at chance in both regimes (\AccRangeR{} on \RJB{}) because
verdicts are dominated by the presentation slot.  After balanced position
calibration, the three smallest judges are significantly \emph{worse} than
chance on \RJB{} (\CalRCoderM{}--\CalRQwenS{}): they prefer the longer patch,
and the defective variant is longer in \VarLongerShare{} of pairs.  Only the
4B judges retain a weak correct signal (up to \CalRQwenM{}).  Accuracy does
not vary measurably with diff size, retrieved repository context and thinking
mode only partly counteract the length preference, and a static-analysis
baseline with eight targeted rules reaches \StaticRelAcc{}.  A trivial
shorter-patch rule scores \LengthHeurAcc{}, which exposes a length artefact
that future versions of the benchmark must remove.  We release the dataset,
harness, prompts and all verdict logs.
